\section{Limitations and Future Work}
\label{sec:limitations}

\textbf{Reactive state resets on topology change.} The circuit for an entire network is
rebuilt from scratch whenever its wiring changes anywhere in that network, which resets a
capacitor's stored charge or an inductor's current to zero at that instant. The memristor is
deliberately exempted, since persistent state is its entire pedagogical point, but a student
who, e.g., extends a circuit while a capacitor is mid-charge will see that capacitor's
voltage reset rather than continue from where it was. Carrying reactive state across a
topology rebuild the same way the memristor's state is carried is planned future work and is
mostly a bookkeeping problem: the new topology's node indices must be matched back to the
element's previous node indices before its stored voltage or current can be meaningfully
re-applied.

\textbf{Fixed presets rather than continuous values.} Every adjustable quantity -- a
resistor's resistance, a source's amplitude or frequency -- is one of a small number of
discrete presets cycled by right-clicking, rather than a numeric entry. This keeps the
interaction model consistent with the rest of the base game (no custom text-entry GUI is
required) at the cost of the fine-grained control a real bench instrument offers; a
numeric-entry interface for these values is a natural extension.

\textbf{Solver scalability.} The MNA system for each independent network is solved by dense
Gaussian elimination every tick, which is adequate for the circuit sizes a single player
typically builds by hand, but scales as $O(n^3)$ in the number of independent nodes and does
not exploit the sparsity that a large circuit's conductance matrix would actually have. A
sparse solver, or amortizing re-solves across ticks when nothing electrically relevant has
changed, would be necessary before the tool could comfortably support very large
student-built circuits or a busy multiplayer server hosting many networks at once.

\textbf{Verification is currently informal.} As noted in Section~\ref{sec:validation}, the
closed-form comparisons used during development have not yet been captured as a persisted,
automatically re-run regression suite, and the tool's pedagogical effectiveness has so far
been assessed only informally by its author, not through a controlled classroom study. A
natural next step is a pre/post evaluation using an established circuits concept inventory,
comparing a cohort using Mine Memristors as a supplementary lab activity against a cohort
using a conventional schematic-based simulator, to test whether the game-based,
first-person construction experience described in Section~\ref{sec:pedagogy} produces a
measurable difference in conceptual understanding rather than only in engagement.

\textbf{Platform dependency.} The mod is necessarily tied to a specific version of
Minecraft and the Fabric mod loader, both of which are externally maintained and continue
to evolve; Section~\ref{sec:validation} discusses the verification discipline this
requires, but ongoing maintenance to track future game updates is an unavoidable cost of
building instructional software on top of an actively-developed commercial platform rather
than a purpose-built, stable simulator.
